\section{Experimental Setup}\label{sec:setup}

\subsection{Design principle: replicate the protocol, not the pixels}

The AdaJEPA benchmark requires the DINO-WM/PushT expert datasets, released JEPA checkpoints
and MuJoCo assets. To make a full study reproducible on a single CPU-only machine we
re-implemented the environments in numpy---an impulse-based 2D rigid-contact simulator
(circle pusher against a polyomino block whose convex decomposition is the set of its unit
cells) and a point-mass maze with grid walls---together with a software rasteriser producing
$64{\times}64$ RGB frames. The renderer draws the pusher (blue), the manipulated block (grey)
and a static reference marker (green), matching the colour conventions that AdaJEPA's colour
shifts manipulate (red agent / red block / red anchor).

Everything that defines the \emph{evaluation protocol} is kept identical to AdaJEPA:
\begin{itemize}[leftmargin=1.2em,itemsep=2pt]
\item goals are \emph{future frames of the same held-out trajectory}, 15 steps ahead in our
runs (AdaJEPA uses 25 with a 100-action budget; we keep a 2:1 executor-to-goal slack instead
of 4:1 because our action cost is fixed, see \S\ref{sec:setup}), with trivial segments
filtered out;
\item MPC uses horizon $H=8$, executes a chunk of $c=5$ actions, at most $R=6$ replans,
and warm-starts the remaining action sequence from the previous plan;
\item the gradient-descent planner optimises the action sequence with Adam at learning rate
$0.1$, initialised at zero;
\item success is a pose/position tolerance in the environment at the end of each MPC step, and
reported as a success-versus-replan curve as well as a final rate;
\item adaptation happens once per replanning step, on a replay buffer of the 5 most recent
transitions, with learning rates $5\times10^{-4}$ (predictor) and $10^{-5}$ (encoder head).
\end{itemize}

\subsection{Environment suites and shift generators}

\begin{table}[t]
\centering
\caption{Shift families, their generator parameters, and the training/evaluation splits. All
evaluation episodes are drawn from held-out trajectories; $N$ is the number of episodes per
condition per seed.}
\label{tab:suites}
\begin{tabular}{llll}
\toprule
Family & Shift knob & Train & Test \\
\midrule
Shape & block geometry & \texttt{T,L,Z,+} (4$\times$350 traj.) & 7 shapes, 3 unseen (\texttt{I}, \texttt{smallT}, \texttt{square})\\
Visual (corruption) & per-frame corruption & \texttt{T} & blur / salt-and-pepper / dark\\
Visual (colour) & entity colour & \texttt{T} & red agent / red block / red anchor\\
Dynamics & mass $\times0.2$, damping $\times20$ & default dynamics, 25 layouts & low mass / high damping\\
Layout & maze connectivity & 25 random $8{\times}8$ layouts & 5 held-out layouts, BFS goals at distance 3--5\\
Continual & environment sequence & --- & $\texttt{T}\!\to\!\texttt{L}\!\to\!\texttt{Z}\!\to\!\texttt{+}\!\to\!\texttt{T}$ (model persists)\\
\bottomrule
\end{tabular}
\end{table}

Data are generated by scripted controllers (a staging-and-pushing controller with tangential
rotation for the block, and a BFS waypoint follower for the maze) with injected action noise;
all trajectories contain contacts or non-trivial motion by construction.

\subsection{World model and training}

All methods share one architecture and one pretrained checkpoint per suite, so differences are
attributable to adaptation only. The model has $\sim$1.3M parameters ($d=128$, two S-blocks,
one R-block, rank-4 LoRA, 594k in the fast set). Pretraining uses Adam ($5\times10^{-4}$,
cosine decay) for 6 epochs on PushBlock windows and 8 epochs on maze windows with batch 48,
$\lambda_{\mathrm{unc}}=0.1$, $\lambda_{\mathrm{reg}}=1.0$, $\lambda_{\mathrm{inv}}=0.1$,
$J=4$ future steps and $K_{\max}=4$ refinement depths. Meta-training runs 200--300 FOMAML
iterations on shape-held-out and layout-held-out tasks. Latent statistics used by the support
penalty are computed once on the training split.

\subsection{Compared methods}

\begin{itemize}[leftmargin=1.2em,itemsep=2pt]
\item \textbf{Frozen}: no adaptation (AdaJEPA's frozen baseline).
\item \textbf{AdaJEPA}: our re-implementation of the paper's default setting---one gradient
step per replan, learning rates $5\times10^{-4}$/$10^{-5}$, direct updates of the predictor's
last transformer block and the encoder head, recent-5 buffer, per-episode reset, uncertainty
terms switched off in the planning cost.
\item \textbf{DTA-JEPA}: full method (dual timescale, uncertainty allocation, dual memory,
safety layer, adaptive refinement depth, cross-episode consolidation).
\item \textbf{Ablations}: \texttt{lora} (rank-4 LoRA fast set instead of direct updates),
\texttt{nounc} (no uncertainty allocation and no uncertainty terms in the cost),
\texttt{nomem} (no episodic or parametric memory), \texttt{nosafe} (no anchor ball, no
rollback, no support penalty), \texttt{fixedk} (fixed refinement depth), \texttt{noslow}
(no cross-episode consolidation or skill bank) and \texttt{nometa} (no FOMAML initialisation).
\end{itemize}

Ablations use the same world model, planner, episode set and budget as the full method, so
each row isolates one mechanism.

\subsection{Metrics}

Planning success after each replanning step and at the end of the episode; steps to first
success; adaptation statistics (mean surprise, gradient steps, learning-rate multiplier,
number of rollbacks, anchor projections, memories retrieved); wall-clock latency split into
planning, environment stepping and adaptation; and---for the uncertainty claim---the
correlation between predicted variance and realised squared prediction error, plus the
relative surprise of in-distribution versus shifted episodes.
